% GENERATED FILE. Edit canonical lessons, then run npm run generate:books.
% canonical-source-hash: fnv1a64:2843be05a7c3bd82
% canonical-lessons: KA-C481-karavali, KA-C481-kolli, KA-C481-havala, KA-C481-cippu, KA-C481-nadimukha

\chapter{Along the Coast}
\label{ch:ka-along-the-coast}
\hlchaptermodality{481}

\begin{chapteropening}
\textbf{By the end of this chapter:} \emph{I can say a coastline, a bay, coral, a shell and an estuary.}

Complete the last lesson of chapter 481: I can say a coastline, a bay, coral, a shell and an estuary.
\end{chapteropening}

\section[\texorpdfstring{\kn{ಕರಾವಳಿ} (karāvaḷi)}{ಕರಾವಳಿ (karāvaḷi)}]{\kn{ಕರಾವಳಿ} (karāvaḷi) — a coastline; a coast}

\label{lesson:KA-C481-karavali}



\begin{quote}
Before the new one: say the Kannada for a plateau, then the Kannada for to wear away.
\end{quote}

\subsection*{You'll want to know: \kn{ಕರಾವಳಿ}}
\textbf{\kn{ಕರಾವಳಿ}} — \emph{karāvaḷi} — \textquotedblleft{}a coastline; a coast\textquotedblright{}.

\begin{quote}
\kn{ಕರಾವಳಿಯಲ್ಲಿ} \kn{ಮೀನುಗಾರರ} \kn{ಹಳ್ಳಿಗಳು} \kn{ಇವೆ}.
\end{quote}

\begin{quote}
\emph{karāvaḷiyalli mīnugārara haḷḷigaḷu ive.}
\end{quote}

\begin{quote}
There are fishing villages along the coast.
\end{quote}

\kn{ಕರಾವಳಿ} \kn{ಪ್ರದೇಶ}, the coastal region.

\begin{cousinweb}
\textbf{\kn{ಕರಾವಳಿ}} (\emph{karāvaḷi}) is a Kannada word for a coastline whose origin is not settled.
\end{cousinweb}

\begin{grammarlens}[title={What you've built}]
One more word for getting around.
\end{grammarlens}

\subsection*{Your turn}
\begin{itemize}
\raggedright
  \item \emph{Say it:} \emph{karāvaḷi}
  \item \emph{Say it:} \emph{karāvaḷi}, once more
  \item \emph{Say it:} \emph{karāvaḷiyalli mīnugārara haḷḷigaḷu ive.} — There are fishing villages along the coast.
  \item \emph{Recall:} say the Kannada for a plateau, then the Kannada for to wear away, then say \emph{karāvaḷi} again
\end{itemize}

\begin{tcolorbox}[breakable,colback=teal!4,colframe=teal!35!black,title={Before you move on}]
What does \kn{ಕರಾವಳಿ} mean? (\textquotedblleft{}A coastline; a coast\textquotedblright{}.) Say it once more.
\end{tcolorbox}
\section[\texorpdfstring{\kn{ಕೊಲ್ಲಿ} (kolli)}{ಕೊಲ್ಲಿ (kolli)}]{\kn{ಕೊಲ್ಲಿ} (kolli) — a bay; a gulf}

\label{lesson:KA-C481-kolli}



\begin{quote}
Before the new one: say the Kannada for to wear away, then the Kannada for a coastline.
\end{quote}

\subsection*{You'll want to know: \kn{ಕೊಲ್ಲಿ}}
\textbf{\kn{ಕೊಲ್ಲಿ}} — \emph{kolli} — \textquotedblleft{}a bay; a gulf\textquotedblright{}.

\begin{quote}
\kn{ಕೊಲ್ಲಿಯಲ್ಲಿ} \kn{ಅಲೆಗಳು} \kn{ಶಾಂತವಾಗಿರುತ್ತವೆ}.
\end{quote}

\begin{quote}
\emph{kolliyalli alegaḷu śāntavāgiruttave.}
\end{quote}

\begin{quote}
The waves are calm in the bay.
\end{quote}

Water held in a curve of the coast.

\begin{cousinweb}
\textbf{\kn{ಕೊಲ್ಲಿ}} (\emph{kolli}) is native Dravidian: DEDR 2137 lists Kannada \emph{kolli}, \emph{kolle}, \textquotedblleft{}a bend, a corner, a gulf, a bay\textquotedblright{}.
\end{cousinweb}

\begin{grammarlens}[title={What you've built}]
One more word for getting around.
\end{grammarlens}

\subsection*{Your turn}
\begin{itemize}
\raggedright
  \item \emph{Say it:} \emph{kolli}
  \item \emph{Say it:} \emph{kolli}, once more
  \item \emph{Say it:} \emph{kolliyalli alegaḷu śāntavāgiruttave.} — The waves are calm in the bay.
  \item \emph{Recall:} say the Kannada for to wear away, then the Kannada for a coastline, then say \emph{kolli} again
\end{itemize}

\begin{tcolorbox}[breakable,colback=teal!4,colframe=teal!35!black,title={Before you move on}]
What does \kn{ಕೊಲ್ಲಿ} mean? (\textquotedblleft{}A bay; a gulf\textquotedblright{}.) Say it once more.
\end{tcolorbox}
\section[\texorpdfstring{\kn{ಹವಳ} (havaḷa)}{ಹವಳ (havaḷa)}]{\kn{ಹವಳ} (havaḷa) — coral}

\label{lesson:KA-C481-havala}



\begin{quote}
Before the new one: say the Kannada for a coastline, then the Kannada for a bay.
\end{quote}

\subsection*{You'll want to know: \kn{ಹವಳ}}
\textbf{\kn{ಹವಳ}} — \emph{havaḷa} — \textquotedblleft{}coral\textquotedblright{}.

\begin{quote}
\kn{ಸಮುದ್ರದ} \kn{ತಳದಲ್ಲಿ} \kn{ಬಣ್ಣಬಣ್ಣದ} \kn{ಹವಳ} \kn{ಇದೆ}.
\end{quote}

\begin{quote}
\emph{samudrada taḷadalli baṇṇabaṇṇada havaḷa ide.}
\end{quote}

\begin{quote}
There is colourful coral on the sea bed.
\end{quote}

\kn{ಹವಳದ} \kn{ಸರ}, a coral necklace.

\begin{cousinweb}
\textbf{\kn{ಹವಳ}} (\emph{havaḷa}) is native Dravidian: DEDR 3998 lists Kannada \emph{pavaḷa}, \textquotedblleft{}coral\textquotedblright{}, with the old \emph{p-} that modern Kannada turns into \emph{h-}. Sanskrit \emph{pravāla}, \textquotedblleft{}coral\textquotedblright{}, was itself taken from Dravidian, as Turner notes (CDIAL 8794).
\end{cousinweb}

\begin{grammarlens}[title={What you've built}]
One more word for getting around.
\end{grammarlens}

\subsection*{Your turn}
\begin{itemize}
\raggedright
  \item \emph{Say it:} \emph{havaḷa}
  \item \emph{Say it:} \emph{havaḷa}, once more
  \item \emph{Say it:} \emph{samudrada taḷadalli baṇṇabaṇṇada havaḷa ide.} — There is colourful coral on the sea bed.
  \item \emph{Recall:} say the Kannada for a coastline, then the Kannada for a bay, then say \emph{havaḷa} again
\end{itemize}

\begin{tcolorbox}[breakable,colback=teal!4,colframe=teal!35!black,title={Before you move on}]
What does \kn{ಹವಳ} mean? (\textquotedblleft{}Coral\textquotedblright{}.) Say it once more.
\end{tcolorbox}
\section[\texorpdfstring{\kn{ಚಿಪ್ಪು} (cippu)}{ಚಿಪ್ಪು (cippu)}]{\kn{ಚಿಪ್ಪು} (cippu) — a shell}

\label{lesson:KA-C481-cippu}



\begin{quote}
Before the new one: say the Kannada for a bay, then the Kannada for coral.
\end{quote}

\subsection*{You'll want to know: \kn{ಚಿಪ್ಪು}}
\textbf{\kn{ಚಿಪ್ಪು}} — \emph{cippu} — \textquotedblleft{}a shell\textquotedblright{}.

\begin{quote}
\kn{ಮಕ್ಕಳು} \kn{ಕಡಲತೀರದಲ್ಲಿ} \kn{ಚಿಪ್ಪುಗಳನ್ನು} \kn{ಸಂಗ್ರಹಿಸಿದರು}.
\end{quote}

\begin{quote}
\emph{makkaḷu kaḍalatīradalli cippugaḷannu saṅgrahisidaru.}
\end{quote}

\begin{quote}
The children collected shells on the beach.
\end{quote}

Plural: \kn{ಚಿಪ್ಪುಗಳು}, \emph{cippugaḷu}.

Also the shell of a coconut: \kn{ತೆಂಗಿನ} \kn{ಚಿಪ್ಪು}.

\begin{cousinweb}
\textbf{\kn{ಚಿಪ್ಪು}} (\emph{cippu}) is native Dravidian: DEDR 2535 lists Kannada \emph{cippu}, \emph{sippu}, \textquotedblleft{}an oyster shell, a mussel, a piece of coconut shell\textquotedblright{}.
\end{cousinweb}

\begin{grammarlens}[title={What you've built}]
One more word for getting around.
\end{grammarlens}

\subsection*{Your turn}
\begin{itemize}
\raggedright
  \item \emph{Say it:} \emph{cippu}
  \item \emph{Say it:} \emph{cippu}, once more
  \item \emph{Say it:} \emph{makkaḷu kaḍalatīradalli cippugaḷannu saṅgrahisidaru.} — The children collected shells on the beach.
  \item \emph{Recall:} say the Kannada for a bay, then the Kannada for coral, then say \emph{cippu} again
\end{itemize}

\begin{tcolorbox}[breakable,colback=teal!4,colframe=teal!35!black,title={Before you move on}]
What does \kn{ಚಿಪ್ಪು} mean? (\textquotedblleft{}A shell\textquotedblright{}.) Say it once more.
\end{tcolorbox}
\section[\texorpdfstring{\kn{ನದೀಮುಖ} (nadīmukha)}{ನದೀಮುಖ (nadīmukha)}]{\kn{ನದೀಮುಖ} (nadīmukha) — an estuary; a river mouth}

\label{lesson:KA-C481-nadimukha}



\begin{quote}
Before the new one: say the Kannada for coral, then the Kannada for a shell.
\end{quote}

\subsection*{You'll want to know: \kn{ನದೀಮುಖ}}
\textbf{\kn{ನದೀಮುಖ}} — \emph{nadīmukha} — \textquotedblleft{}an estuary; a river mouth\textquotedblright{}.

\begin{quote}
\kn{ನದೀಮುಖದಲ್ಲಿ} \kn{ನದಿ} \kn{ಸಮುದ್ರವನ್ನು} \kn{ಸೇರುತ್ತದೆ}.
\end{quote}

\begin{quote}
\emph{nadīmukhadalli nadi samudravannu sēruttade.}
\end{quote}

\begin{quote}
At the estuary the river joins the sea.
\end{quote}

\kn{ನದಿ}, a river, plus \kn{ಮುಖ}, a mouth.

\begin{cousinweb}
\textbf{\kn{ನದೀಮುಖ}} (\emph{nadīmukha}) is built from two Sanskrit words: \emph{nadī}, \textquotedblleft{}a river\textquotedblright{} (CDIAL 6943), and \emph{mukha}, \textquotedblleft{}a mouth, an entrance\textquotedblright{} (CDIAL 10158).
\end{cousinweb}

\begin{grammarlens}[title={What you've built}]
One more word for getting around.
\end{grammarlens}

\subsection*{Your turn}
\begin{itemize}
\raggedright
  \item \emph{Say it:} \emph{nadīmukha}
  \item \emph{Say it:} \emph{nadīmukha}, once more
  \item \emph{Say it:} \emph{nadīmukhadalli nadi samudravannu sēruttade.} — At the estuary the river joins the sea.
  \item \emph{Recall:} say the Kannada for coral, then the Kannada for a shell, then say \emph{nadīmukha} again
\end{itemize}

\begin{tcolorbox}[breakable,colback=teal!4,colframe=teal!35!black,title={Before you move on}]
What does \kn{ನದೀಮುಖ} mean? (\textquotedblleft{}An estuary; a river mouth\textquotedblright{}.) Say it once more.
\end{tcolorbox}
